\section{The original class number and the column-rank producer}
\label{sec:rank-producer}

\status{RMP1--RMP11 prove the literal \texttt{ClassNumberFormula}$\,{}_2$
for every supplied original $D:\texttt{D7}$ and
$\mathrm{Lv}:\texttt{LevelData}(D)$. They then prove the original
\texttt{RankM} from these records and level-zero column forms on the
same $D$. The source's already supplied \texttt{TSel}, or its original
\texttt{RawSel Bundle}, supplies those forms with one common threshold.
For the selected construction class, its membership predicate itself
supplies this Bundle through \texttt{cls\_at} and \texttt{pick}.
Thus the separate rank-formula producer is discharged there without
a further selector premise. This does not construct a member of the
common arithmetic family class.
The mass input is the previously reconstructed MF18--MF19; the genus
dictionary is GV11. No new mass theorem or Tamagawa theorem is invoked.}

\paragraph{Notation and exact target.}
Fix the original datum $D$ and a word $u$. Write
\[
 \begin{gathered}
 E=D.\mathrm{host}.E,\quad F=D.\mathrm{host}.F,\quad
 k=D.R.k,\quad f=D.R.f=[F:\mathbb Q],\\
 T=(D.\mathrm{mem}(u)).A.T,\qquad g_i=D.\mathrm{host}.g(i).
 \end{gathered}
\]
The Hermitian form is conjugate linear in its first variable.
In the original coordinates $\mathrm{crd}_0$ it has Gram matrix
$J=D.\mathrm{sp}.\mathrm{gram}_0$. Put
\[
 \begin{gathered}
 G=\mathrm U(J)(\mathbb A_{F,f}),\qquad
 \Gamma=\mathrm U(J)(F),\qquad L=\mathrm{Lv}(u).\mathrm{Lb},\\
 U=\mathrm{Lv}(u).U,\qquad C_u=(D.\mathrm{mem}(u)).C.\mathrm{cls}(0).
 \end{gathered}
                                                               \tag{RMP.1}
\]
Here $\mathrm U(J)(F)$ means the rational unitary matrices with entries
in $E$; it is the source's \texttt{ratPts} inside \texttt{UPts}.
The notation $L^{\#}$ denotes the actual Hermitian dual
\texttt{dualSet}, which is already an $\mathcal O_E$-module and
therefore equals its span \texttt{dualLb}.
The target class number is
\[
 |C_u|=\mathrm{massI}(D,k)\,
       \mathrm{idxPi}(D,g_0,g_1)(T^k-1)^f.              \tag{RMP.2}
\]
The rank target is the literal definition
\[
 \mathrm{RankM}(D):\quad
 D.\mathrm{duS}(u)=\mathrm{duStarU}(D,u)(T^k-1)^f
       \quad\hbox{for every }u,                        \tag{RMP.3}
\]
where $D.\mathrm{duS}(u)=\dim_{\mathbb C}(D.\mathrm{mem}(u)).A.\mathrm{Mu}_0$.
All equalities involving the source's real-valued factors use the
natural embedding of the integer cardinalities into $\mathbb R$.

% BEGIN RMP1
\paragraph{RMP1: the exact base, primes and level conditions.}
\label{claim:RMP1}
The original $\mathrm{ColLevel}$ gives a finitely generated full lattice $L$ and
\[
 L\subset L^{\#},\qquad T L^{\#}\subset L,\qquad
 |L^{\#}/L|=T^{2f}.                                    \tag{RMP.4}
\]
It gives the local groups $U_v$, with $U_v\subset\operatorname{Stab}(L_v)$,
and identifies $U$ with the unitary adeles having components in all
$U_v$. Outside the primes above $g_0,g_1,T$ it gives
$U_v=\operatorname{Stab}(L_v)$. Above $g_i$, an actual split projection
and an actual frame identify $U_v$ with the inverse image of the
upper-unitriangular group under reduction
$\mathrm{GL}_k(\mathcal O_{F_v})\to\mathrm{GL}_k(k_v)$.
Above $T$ it gives an $\mathcal O_{E,v}$-module $M_v$ of additive
cardinality $T$ and a nonzero linear functional
\[
 t_v:L_v^{\#}\longrightarrow M_v,\qquad
 t_v(L_v)=0,\qquad t_v(TL_v^{\#})=0,
\]
and, literally,
\[
 U_v=\operatorname{Stab}(L_v)\cap\operatorname{Stab}(t_v).
                                                               \tag{RMP.5}
\]
The last stabilizer fixes the functional itself, not just its kernel.
These are stored fields, not proposed replacements for the level.
[Sources R01--R03.]

The source's $g_0,g_1$ are distinct rational primes. They split
completely in $E$: the field \texttt{nd.g\_split} says that the
number of primes is $[E:\mathbb Q]$. The degree formula
$\sum_{P\mid g_i}e_P f_P=[E:\mathbb Q]$, with positive integral
summands, forces every $e_P=f_P=1$. The ramification and residue-degree
multiplication formulas show the same in $F$; hence there are exactly
$f$ places of $F$ over $g_i$, each of residue cardinality $g_i$.

For $T$ the source stores $\pi\bar\pi=T$ and
$\pi\equiv1\pmod{N\mathcal O_E}$, where
$N=D.\mathrm{host}.\mathrm{QS}\,q\,d\,(g_0g_1)^2$.
Conjugation preserves this congruence, so $T-1\in N\mathcal O_E$.
A rational algebraic integer is an integer; thus $N\mid T-1$ in
$\mathbb Z$. In particular $T\ne g_0,g_1$ and
$T\equiv1\pmod{\mathrm{condN}}$.
The original host's Frobenius/Artin fields identify the Frobenius at
this unramified prime with the Artin image of its residue class.
That class is $1$, so Frobenius is $1$; the stored unramifiedness
makes inertia trivial. The decomposition group, generated modulo
inertia by Frobenius, is trivial. Thus $T$ splits completely in $E$
and in $F$, proving the source's \texttt{SplitsF} statements without
a prime-existence assertion. Equivalently, the later argument can
retain exactly the two \texttt{SplitsF} hypotheses of (RMP.2).
[Sources R04--R06.]

Finally, the base space $D.\mathrm{sp}.X(0)$ is totally positive
definite, not just definite of an unspecified sign. Its rank is $k$;
its supplied definite signature is either $(k,0)$ or $(0,k)$.
It embeds isometrically into $Y(0)$, whose signature is
$(k+r_S,r_S+1)$ with $r_S+1=m/2<k$. A negative-definite subspace of
rank $k$ would exceed that negative index. This proves the positive
alternative at every embedding by Sylvester's law of inertia.
The source's unramified CM condition is exactly the condition treated
in MF1 and DR1, including the dyadic places. [Sources R02, R07.]
% END RMP1

% BEGIN RMP2
\paragraph{RMP2: the actual local discriminant module has length two.}
\label{claim:RMP2}
Fix $v\mid T$, put $R=\mathcal O_{F_v}$, let $\varpi$ be its
uniformizer and let $k_v=\mathbb F_T$. Write
$B=\mathcal O_E\otimes_{\mathcal O_F}R$.
The unramified quadratic algebra $B$ is either $R\times R$ or an
unramified quadratic DVR. The action on the cyclic additive group
$M_v$ of prime order defines a unital ring map
$B\to\operatorname{End}_{\mathbb Z}(M_v)=\mathbb F_T$.
Its restriction to $R$ is reduction: its kernel is the unique maximal
ideal. An unramified quadratic residue field cannot embed in
$\mathbb F_T$. Thus $B=R\times R$, and the action on $M_v$ selects
one factor. This gives the exact ordering of the split factors
used by the source's \texttt{exists\_local}; it is not an arbitrary
choice subsequently imposed on $t_v$.

Let $e_+,e_-$ be the corresponding integral idempotents, with
$e_+M_v=M_v$ and $e_-M_v=0$. Under the two projections write
\[
 L_v=A_v\oplus B_v,\qquad
 L_v^{\#}=B_v^{\vee}\oplus A_v^{\vee}.                 \tag{RMP.6}
\]
The superscript $\vee$ uses the nondegenerate cross-pairing given
by $J$. The two same-factor pairings vanish because conjugation
exchanges $e_+$ and $e_-$. The two projected modules are free
$R$-lattices of rank $k$, and the two quotient lengths
\[
 d_v=\operatorname{length}_R(B_v^{\vee}/A_v)
     =\operatorname{length}_R(A_v^{\vee}/B_v)           \tag{RMP.7}
\]
are equal. To check this directly, choose a basis of $B_v$, use its
dual basis for $B_v^{\vee}$, and diagonalize the full sublattice
$A_v$ by elementary divisors. If the diagonal exponents are
$a_1,\ldots,a_k\ge0$, its dual enlarges $B_v$ by
$\varpi^{-a_1},\ldots,\varpi^{-a_k}$. Both lengths are
$\sum_j a_j$. This also proves that duality commutes with the
localization used in (RMP.6).

Because $t_v\ne0$, its image is all of the prime-order module
$M_v$. It factors through the first quotient in (RMP.6), so $d_v\ge1$.
Away from $T$, (RMP.4) makes $L_v^{\#}=L_v$. Localization and the
Chinese remainder theorem for the finite module $L^{\#}/L$ give
\[
 T^{2f}=|L^{\#}/L|=\prod_{v\mid T}T^{2d_v}.
                                                               \tag{RMP.8}
\]
There are exactly $f$ such places, and all $d_v\ge1$. Taking
integer exponents in (RMP.8) gives $d_v=1$ at every one of them.
Thus $t_v$ induces an isomorphism
$B_v^{\vee}/A_v\simeq M_v$ and kills the entire second component.
In particular its first-component kernel is exactly $A_v$.
The two length-one quotients also give the exact $B$-linear
discriminant-module isomorphism to the product of the two simple
residue modules in \texttt{LocalAtT}.
This is the kernel equality for which the source uses its
\texttt{mem\_of\_count} argument. [Sources R01, R08--R10.]
% END RMP2

% BEGIN RMP3
\paragraph{RMP3: construct the self-dual enlargement on the original base.}
\label{claim:RMP3}
At $v\mid T$ set
\[
 N_v=B_v^{\vee}\oplus B_v
     =e_+L_v^{\#}+e_-L_v;
 \qquad N_v=L_v\quad(v\nmid T).                       \tag{RMP.9}
\]
These local lattices are self-dual by the cross-pairing in (RMP.6).
Define the global lattice by the literal intersection formula
\[
 N=\{x\in L^{\#}:e_-x_v\in L_v\text{ for every }v\mid T\}.
                                                               \tag{RMP.10}
\]
This is exactly \texttt{overLat} with the chosen negative-factor
idempotents. It contains $L$ and lies in $L^{\#}$. Multiplication by
$T$ embeds $L^{\#}$ into the finitely generated $\mathcal O_E$-module $L$;
since $\mathcal O_E$ is Noetherian, $L^{\#}$ and its submodule $N$
are finitely generated. Their spans are full because
they contain $L$.

The completion of (RMP.10) is precisely (RMP.9). One way to see both
inclusions without any approximation assumption is to use
$Q=L^{\#}/L$. It is a finite module supported above $T$.
Its canonical Chinese-remainder decomposition is the direct sum of
its localizations $Q_v$. Choose in each $Q_v$ the submodule
$e_+Q_v$, and take the inverse image of their direct sum in
$L^{\#}$. This inverse image is (RMP.10), and its localizations
are (RMP.9); completing these finite-length quotients changes
nothing. This is also the finite-lattice patching argument of LP1.

Membership of an element of $E$ in $\mathcal O_E$ is tested at every
finite place. Applying that fact to all pairings, and using
$N_v^{\#}=N_v$, gives $N^{\#}=N$. Thus $N$ is an actual full
self-dual lattice on the original Hermitian space.

At $v\mid T$, a unitary map preserving $L_v$ preserves its dual
and its two factors, hence preserves $N_v$. Conversely a unitary
map preserving $N_v$ is determined by its action on $B_v^{\vee}$,
with the adjoint inverse action on $B_v$. Let
\[
 f_v:B_v^{\vee}\longrightarrow k_v
                                                               \tag{RMP.11}
\]
be the original $t_v$ transported through any additive
$M_v\simeq k_v$ respecting the first-factor action. Its kernel is
$A_v$. Fixing $f_v$ therefore preserves $A_v$, and duality also
preserves $A_v^{\vee}$. Consequently (RMP.5) is exactly
\[
 U_v=\{a\in\operatorname{Stab}(N_v):f_v\circ a_+=f_v\}.
                                                               \tag{RMP.12}
\]
Away from $T$, $N_v=L_v$ and the original inclusion
$U_v\subset\operatorname{Stab}(L_v)$ suffices. Hence
$U\subset K_N:=\operatorname{Stab}(\widehat N)$.
All assertions of \texttt{IsOverLattice}, the projection and
functional clauses of \texttt{LocalAtT}, and the inclusions in
\texttt{HasOverLattice} now hold on this very $N$.
The unique unitary lifting of a first-factor matrix is given by
its adjoint inverse on the second factor; a basis of
$B_v^{\vee}$ gives the required frame. [Sources R08--R11.]
% END RMP3

% BEGIN RMP4
\paragraph{RMP4: the local index above $T$ is $T^k-1$.}
\label{claim:RMP4}
In a basis of $B_v^{\vee}$ the first split projection identifies
$K_{N,v}$ with $\mathrm{GL}_k(R)$. It identifies $f_v$ with a
nonzero row $r_v\in k_v^k$. Indeed $f_v$ is surjective and its
kernel has quotient of cardinality $T$, so it factors through
reduction and remains nonzero there. Equation (RMP.12) identifies
$U_v$ with the inverse image of
\[
 \{a\in\mathrm{GL}_k(k_v):r_va=r_v\}.                \tag{RMP.13}
\]
The reduction map from $\mathrm{GL}_k(R)$ is onto: lift any
invertible residue matrix entrywise; its determinant is a unit.
The group $\mathrm{GL}_k(k_v)$ acts transitively on the nonzero
rows, since any such row can be extended to a dual basis. Its
orbit therefore has $T^k-1$ elements. The ordinary orbit-stabilizer
bijection now gives
\[
                       [K_{N,v}:U_v]=T^k-1.           \tag{RMP.14}
\]
There is no division by $T-1$: the level fixes a specified nonzero
functional, not its projective line. This proves the exact
\texttt{IndexAtT} and local-index conclusion used in the original
class-number calculation. [Sources R11--R12.]
% END RMP4

% BEGIN RMP5
\paragraph{RMP5: auxiliary indices and their exact adelic product.}
\label{claim:RMP5}
At a place above $g_i$, the actual split projection and frame give
$U_v$ as the inverse image of the upper-unitriangular subgroup
$N_k(\mathbb F_{g_i})$. This is a subgroup: multiplication preserves
zero entries below the diagonal and diagonal entries equal to one,
and the inverse of $1+X$ is the finite series
$1-X+\cdots+(-X)^{k-1}$ for strictly upper-triangular $X$.
It has $g_i^{k(k-1)/2}$ elements. Relative to the maximal compact
in that frame, reduction therefore gives the index
\[
 \frac{|\mathrm{GL}_k(\mathbb F_{g_i})|}
      {|N_k(\mathbb F_{g_i})|}
       =\mathrm{idxColLoc}(g_i,k).                    \tag{RMP.15}
\]
That maximal compact and $K_{N,v}$ need not have been specified by
the same frame. Both contain $U_v$, and they are conjugate in the
local unitary group: the split first component is simply
$\mathrm{GL}_k(F_v)$, and any two full lattice bases are related
by an invertible matrix. GV1 proves this group unimodular, so its
Haar measure is invariant under conjugation. The two compact groups
have equal measure. Since $U_v$ is compact open, dividing by its
positive measure proves that (RMP.15) is also $[K_{N,v}:U_v]$.

Away from $g_0,g_1,T$, (RMP.9) and the original
\texttt{Uloc\_hyp} give $U_v=K_{N,v}$.
The three sets of exceptional places are disjoint by RMP1.
The map from $K_N/U$ to the product of local right-coset sets is
bijective. Equality is tested locally by the definition of $U$.
For surjectivity choose representatives at the finitely many
exceptional places and put the identity elsewhere; that is an
actual finite unitary adele preserving $\widehat N$.
There are $f$ places over each of the three rational primes.
Thus
\[
 \boxed{[K_N:U]=\mathrm{idxPi}(D,g_0,g_1)(T^k-1)^f.}   \tag{RMP.16}
\]
Both factors are positive and the index is finite.

This also proves the source's stronger
\texttt{StabiliserLevelIndexFormula}, which quantifies over
\emph{any} full self-dual $N'$ with $U\subset K_{N'}$.
DR1 and the split dual-basis construction give a local unitary
isometry $N_v\to N'_v$ at every finite place. The lattices agree
outside a finite set, so these isometries assemble to an adele
conjugating $K_N$ to $K_{N'}$. GV1 and GV12 give equal Haar
measures. Dividing by the measure of the same compact open $U$
shows $[K_{N'}:U]=[K_N:U]$ and preserves (RMP.16).
No global integral orthonormal basis of $N$ or $N'$ is assumed.
[Sources R12--R14.]
% END RMP5

% BEGIN RMP6
\paragraph{RMP6: the chosen self-dual genus has exactly the required mass.}
\label{claim:RMP6}
Transport $N$ through the original coordinate isomorphism:
\[
              M=\mathrm{crd}_0^{-1}(N)
                  \subset(D.\mathrm{sp}.X(0)).V.       \tag{RMP.17}
\]
The displayed identity
$h(x,y)=\operatorname{hform}(J,\mathrm{crd}_0x,\mathrm{crd}_0y)$
shows that finite generation, full span and self-duality are
preserved. Thus (RMP.17) is precisely an admissible $M$ in the
source first-space mass identity, with $\mathrm{crdLat}_0(M)=N$. We have not asserted that the original
standard lattice $X(0).L_0$ is fixed by $U$; (RMP.17) supplies the
self-dual lattice for which that property was proved.

MF18 gives, for every positive definite rank-$k$ integral lattice
in this same unramified CM extension,
$\operatorname{mass}(M)=c_k\lambda(M)$.
Self-duality makes the local discriminant length $d_M$ zero at
every place. MF6 and MF14 then give
$\lambda_v(M)=1$, including the dyadic places; hence
$\lambda(M)=1$. The actual standard coordinate form and lattice
in the definition of $\mathrm{massI}(D,k)$ have the same property.
MF19 identifies their mass with $c_k$. Therefore
\[
             \operatorname{mass}(M)=\mathrm{massI}(D,k).
                                                               \tag{RMP.18}
\]
This argument uses the already proved uniformity in the rational
Hermitian space in MF18. It does not posit a rational or integral
isometry from the supplied positive definite space to the standard
one. Its inherited named classical input is exactly MF4's
Hasse--Minkowski theorem; no additional absolute mass formula is
needed. [Sources R02, R07, R15--R16; inherited MF6, MF14,
MF18--MF19.]
% END RMP6

% BEGIN RMP7
\paragraph{RMP7: the actual double-coset dictionary and finite stabilizers.}
\label{claim:RMP7}
For $a\in G$, intersect the adelic lattice $a\widehat N$ with
$E^k$. LP1 shows that its localizations are the prescribed ones.
The map
\[
 [a]\longmapsto[\,E^k\cap a\widehat N\,]
                                                               \tag{RMP.19}
\]
is the exact bijection from $\Gamma\backslash G/K_N$ to the
genus classes of $N$. Changing $a$ on the right by $K_N$ changes
no lattice; changing it on the left by $\Gamma$ gives a rational
isometry. Conversely equality up to a rational isometry is tested
at every completion and gives the required double-coset equality.
Every genus lattice supplies local isometries, equal to the identity
outside finitely many places, so it supplies an adele. If one uses
the source's rational-local definition of genus, MF3's Cayley
weak-approximation argument meets the relevant open local
transporter and gives the same local lattice. This is the precise
dictionary proved in GV11, not a strong-approximation assertion.

The rational stabilizer
$\Gamma_a=\Gamma\cap aK_Na^{-1}$ is exactly the automorphism
group of $E^k\cap a\widehat N$. It is finite. Indeed this is a
full integer lattice, and the sum over embeddings of its positive
Hermitian norms is a positive-definite real quadratic form.
An automorphism takes each member of an integer basis into a
discrete lattice with its norm fixed. A bounded set contains only
finitely many lattice points, so there are only finitely many
possible images of that basis.
The quotient $\Gamma\backslash G/K_N$ is also finite: it is the
surjective image of the original finite $C_u\simeq\Gamma\backslash G/U$
under enlargement of the right subgroup. Thus the genus mass is
literally the finite sum
\[
 \sum_{[a]\in\Gamma\backslash G/K_N}|\Gamma_a|^{-1}
             =\mathrm{massI}(D,k),                    \tag{RMP.20}
\]
using (RMP.18). The source's fallback convention for mass is not
entered, because actual finite representatives have been supplied.
[Sources R01, R15, R17--R18; inherited GV11, LP1, MF3.]

We also need the smaller level stabilizers to be trivial.
For $b\in G$, the group $\Gamma\cap bUb^{-1}$ preserves the
full lattice $E^k\cap b\widehat L$, so the same argument makes
it finite. At a place above $g_i$, conjugating by the original
frame puts every element into an integral matrix group whose
reduction is upper unitriangular.
If an element there has finite order $g_i^a c$ with
$(c,g_i)=1$, its $g_i^a$th power has prime-to-$g_i$ order and
upper-unipotent reduction. That reduction must be the identity.
Write the power as $1+X$, $X\in\varpi M_k(R)$. If
$X\in\varpi^rM_k(R)$, its $c$th power is
$1+cX$ modulo $\varpi^{r+1}$. Since $c$ is a unit, induction
puts $X$ in every power of $\varpi$, whence $X=0$.
Its order is consequently a power of $g_i$. The projection of a
rational matrix into either split completion is injective, so its
rational order is a power of both distinct $g_0,g_1$ and equals one.
This proves
\[
                    \Gamma\cap bUb^{-1}=\{1\}
                         \quad(b\in G).               \tag{RMP.21}
\]
The original \texttt{stabI\_inj} and \texttt{stabI\_range} fields
identify the stored $C_u$ stabilizers with these groups; hence
$(D.\mathrm{mem}(u)).C.\mathrm{StabZero}$ holds.
This is CR4 on the actual \texttt{ColLevel}. It needs both given
auxiliary primes and definiteness, but no lower bound called
\texttt{GBig}. [Sources R01, R19; inherited CR4 and RSL7.]
% END RMP7

% BEGIN RMP8
\paragraph{RMP8: count the fibers and obtain \texttt{ClassNumberFormula}$\,{}_2$.}
\label{claim:RMP8}
For a representative $a$ of a class in $\Gamma\backslash G/K_N$,
its fiber in $\Gamma\backslash G/U$ is the orbit set of $K_N/U$
under $a^{-1}\Gamma_a a$. The action is free: an element fixing
$kU$ belongs after conjugation to $\Gamma\cap akU(ak)^{-1}$,
which is trivial by (RMP.21). Every orbit has $|\Gamma_a|$
elements, so the fiber cardinality is $[K_N:U]/|\Gamma_a|$.
Summing these finite cardinalities gives
\[
 \begin{split}
 |\Gamma\backslash G/U|
   &=[K_N:U]\sum_{[a]\in\Gamma\backslash G/K_N}|\Gamma_a|^{-1}\\
   &=\mathrm{idxPi}(D,g_0,g_1)(T^k-1)^f\,
                                      \mathrm{massI}(D,k).
 \end{split}                                                     \tag{RMP.22}
\]
The bijection here is the original \texttt{Lv(u).y\_bij} from
$C_u$; its representatives and stabilizer injections were used in
RMP7. Thus (RMP.22) proves exactly (RMP.2).

Introducing the hypotheses of the source definition therefore
proves its literal class-number predicate
$\texttt{ClassNumberFormula}_2(D,\mathrm{Lv})$.
Its initial \texttt{GBig} implication is harmless: the proof
has established a stronger assertion without it. All splitting
hypotheses have either been retained during the local count or
checked from $D$ in RMP1. In particular, for every original family,
every $x>\mathcal F.x_1$ and every actual $i$, this proves the
class-number formula for
$(\mathcal F.D(x,i),(\mathcal F.\mathrm{ex}(x,i)).\mathrm{lev})$.
This is the original \texttt{hn4} slot of \texttt{link\_ex},
with no eventual arithmetic producer left inside that slot.
[Sources R03, R15, R20--R21.]
% END RMP8

% BEGIN RMP9
\paragraph{RMP9: the actual level-zero values have the stated Weyl dimension.}
\label{claim:RMP9}
Supply $\mathcal T_u:\texttt{Forms0}(D,u)$ on this same $D$.
Its \texttt{val0.eY} identifies the original value-variable set
with the original frame space, and \texttt{val0.Fi}, together
with \texttt{Fi\_zero}, identifies the original $\mathrm{Mu}_0$
with the restriction of adelic forms having polynomial coefficient
values of weights
$\lambda_\tau=(D.\mathrm{fm}(u)).\mathrm{lam}(0,\tau)$.
This is an equality of the specified subspaces after that displayed
linear change of variables; a separate representation of the same
abstract dimension would not suffice. [Sources R22--R24.]

RMP7 made the stored stabilizers $C_u.\mathrm{stab}(0,c)$ trivial.
The \texttt{cls0.stabI\_range} field of $\mathcal T_u$ therefore
makes its own rational level stabilizers trivial as well. Its
\texttt{cls0.y\_bij} is indexed by precisely the same $C_u$.
Consequently arbitrary coefficient values at these representatives
extend uniquely to an adelic form. Explicitly, write
$x=\gamma y_c\kappa$ and set
\[
       F(x,z)=\chi(\kappa)\,v_c(\gamma^{-1}z).
                                                               \tag{RMP.23}
\]
The class determines $c$, and the trivial stabilizer makes
$\gamma,\kappa$ unique. The character law and the coefficient
action give the required equivariance; the invertible coefficient
action preserves the polynomial coefficient space. Restriction
and this extension are inverse linear maps. This is the argument
of CR2 and CR5--CR6, now applied to the literal \texttt{Mu0\_iff}
and the literal $C_u$.

At each embedding the coefficient space is the characteristic-zero
Weyl space with dimension
\[
 \prod_{0\le a<b<k}
         \frac{\lambda_\tau(a)-\lambda_\tau(b)+b-a}{b-a}.
                                                               \tag{RMP.24}
\]
CB2--CB3 and IW3/CR1 give this dimension for the actual polynomial
span, including determinant shifts. Restricting its polynomial
functions to invertible frames is injective because those frames
form a Zariski-open dense set; multiplying by the nonzero determinant
power preserves injectivity. Tensoring the coefficient spaces over
$\tau$ multiplies their dimensions.
The original weight equation is
\[
 \lambda_\tau(a)=u_\tau(a)-\left(\frac{k-1}{2}-a\right),
       \qquad u_\tau=\mathrm{uColR}(D,u,\tau).
                                                               \tag{RMP.25}
\]
Thus the numerator in (RMP.24) is exactly $u_\tau(a)-u_\tau(b)$;
its denominator product is $\prod_{j=0}^{k-1}j!$.
These are precisely the source's \texttt{vand} and
\texttt{facProd}. We obtain
\[
 \boxed{\dim_{\mathbb C}\mathrm{Mu}_0
            =|C_u|\prod_\tau\mathrm{dimV}(\mathrm{uColR}(D,u,\tau)).}
                                                               \tag{RMP.26}
\]
This proves the original \texttt{Mu0Tie}. For this dimension
comparison the two supplied level descriptions need only be
indexed by the same stored $C_u$ with its stored stabilizers and
weights; no additional integral or metric equality is used.
The original \texttt{ColsTiedSym} supplies stronger same-carrier
identities when they are needed elsewhere. [Sources R22--R26.]
% END RMP9

% BEGIN RMP10
\paragraph{RMP10: the literal \texttt{RankM} and the stored $d_\star$.}
\label{claim:RMP10}
Supply $\mathrm{Lv}:\texttt{LevelData}(D)$ and
$\mathcal T_u:\texttt{Forms0}(D,u)$ for every word, all on the
original $D$. Multiply (RMP.2) by the product in (RMP.26).
The source defines
\[
 \mathrm{duStarU}(D,u)
   =\mathrm{massI}(D,k)\,\mathrm{idxPi}(D,g_0,g_1)
        \prod_\tau\mathrm{dimV}(\mathrm{uColR}(D,u,\tau)).
                                                               \tag{RMP.27}
\]
Together with $D.\mathrm{duS}(u)=\dim_{\mathbb C}\mathrm{Mu}_0$,
this gives (RMP.3) for every word. This is a pointwise proof of the
original \texttt{RankM}, not an asymptotic estimate and not the
ramified dimension $(D.\mathrm{mem}(u)).A.\mathrm{du}$ of RSL10.

For an original family member with $x>\mathcal F.x_1$, its actual
extra record supplies
\[
  ((D.\mathrm{mem}(u)).A.\mathrm{duStar}:\mathbb R)
       =\mathrm{duStarU}(D,u).                         \tag{RMP.28}
\]
Consequently (RMP.3) is also the literal
$d_M(u)=d_\star(u)(T_u^k-1)^f$ used in DEX5.
No alternative standard lattice, different mass normalization,
or independently chosen number called $d_\star$ has been substituted.
For any supplied $C:\texttt{ColsTiedSym}(D)$, take
$\mathcal T_u=(C.\mathrm{forms}(u)).\mathrm{to0}$.
TG12 supplies this $C$ on its original true-column input with its
stated laws and later-stage record. RMP10 adds no existence demand
beyond those inputs. [Sources R03, R26--R28.]
% END RMP10

% BEGIN RMP11
\paragraph{RMP11: common-selector thresholds and the closed producer.}
\label{claim:RMP11}
Let $\mathcal A$ be the original source class and retain its
specified $s:\texttt{TSel}(\mathcal A)$. For a family
$\mathcal F\in\mathcal A$, write $a_s=(s\mathcal F).1$.
The definition supplies
\[
 R_{x,i}:\texttt{RecT}(\mathcal F.\mathrm{ex}(x,i))
 \quad\text{for all }x>\mathcal F.x_1,\ a_s\le x,
                         \ i\in\mathcal F.I_x.         \tag{RMP.29}
\]
Its actual field $R_{x,i}.T.\mathrm{cols}$ is a
$\texttt{ColsTiedSym}(\mathcal F.D(x,i))$. Use these very columns,
and the level data from that very $\mathcal F.\mathrm{ex}(x,i)$,
in RMP10. Thus with
\[
                   a=\max\{a_s,\mathcal F.x_1+1\}
                                                               \tag{RMP.30}
\]
we have, in the original quantifier order,
\[
 \forall x\ge a\ \forall i\in\mathcal F.I_x\
                      \mathrm{RankM}(\mathcal F.D(x,i)).
                                                               \tag{RMP.31}
\]
The threshold precedes $i$ and every word $u$. No \texttt{Duty},
finite-character assertion, theta square, metric transport, or
additional size hypothesis is needed for this rank consumer.
For \texttt{RawSel Bundle}, the source defines $s$ with threshold
$\mathcal F.x_1$ and selects exactly the Bundle's $R$ field;
then $a=\mathcal F.x_1+1$ suffices. [Sources R29--R31.]

Fix the target condition $\mathcal G=G.\mathrm{cls}(\lambda)$.
The selected design classes are
\[
 \mathcal C^{\prime}=\texttt{WithChN}(\mathcal A\cap\mathcal G),
 \qquad
 \mathcal C_{\mathrm{des}}=\texttt{WithTw}(\texttt{WithSm}
                 (\mathcal C^{\prime})).
\]
They map to $\mathcal A$ by their literal conjunction projections. Apply (RMP.31) to the same family
under this projection. In the definition of \texttt{Closed}, choose
$\varepsilon_1=1$ and the transported bound with all components zero.
For every $0<\varepsilon\le1$, every family and every $C_H$,
class membership permits evaluation of the fixed $s$; the argument works regardless of
its additional \texttt{Large} proof. This proves
\[
 \texttt{Closed}\bigl(\mathcal C_{\mathrm{des}},\
   \mathcal F\mapsto
     \exists a\ \forall x\ge a\ \forall i,
                     \mathrm{RankM}(\mathcal F.D(x,i))\bigr).
                                                               \tag{RMP.32}
\]
The all-member class-number assertion in RMP8 is itself closed
with the same choices, even before $s$ is supplied.
If one instead has only a \emph{closed} eventual Forms0 supplier,
its own $\varepsilon$ and transported bound must be retained;
joining that supplier with the finite implication preserves them
and only replaces the inner threshold by its maximum with
$\mathcal F.x_1+1$. They cannot be reset to zero in that weaker scope.

For the actual construction class there is no additional existence
premise for this selector. Put
$\mathcal A=\texttt{cls}(Q,Lk,E_{\mathrm{row}})$ with
$Bd=\texttt{Bundle}$. A proof $h:\mathcal A(\mathcal F)$
contains a parameter $\kappa$, a \texttt{Link} from $\kappa$ to
$\mathcal F$, and the pointwise \texttt{QOf} property.
The same \texttt{Link} contains $\kappa.\mathrm{Pos}$, so the
positivity guard in \texttt{QOf} is discharged. For each
$x>\mathcal F.x_1$ and $i$, evaluate \texttt{QOf} at
$c=\mathcal F.C_{\mathrm{sz}}$. It supplies a Bundle on the
literal dependent pair
$(\mathcal F.D(x,i),\texttt{objAt}(\mathcal F,x,i))$.
Choose these witnesses pointwise. This is precisely
\texttt{cls\_at} followed by the source's \texttt{pick}, giving
$\sigma:\texttt{RawSel Bundle}(\mathcal A)$.
The \texttt{ClassSel} construction uses this very class and this
very pick (with its actual combined slot in place of $Q$).
The selector of RMP11 is then the \texttt{sel} of $\sigma$:
its selected record is exactly $(\sigma(\mathcal F,h,x,i)).R$.
Consequently its columns, extra record, level groups and stored
$d_\star$ are the ones used throughout RMP8--RMP10.
This verifies (RMP.32) on the selected construction class and its
actual target intersection, without an independent selector premise.

This discharges the v0.31 leaf \texttt{DESIGN\_RANKM\_CLOSED}.
The shared \texttt{DESIGN\_CHOICES\_COMMON} obligation still has
to construct a family satisfying this class together with the
other required common conditions. The argument above proves a
property of every member of the class; it does not assert that
the class has a member. The design-choice and level-size tasks
keep their separate scopes. There is no remaining class-number,
local-index or Weyl-dimension obligation hidden in this rank leaf.
[Sources R20--R21, R29--R33; inherited DEX5--DEX7 and TG12.]
% END RMP11
